\documentclass{ip-lab}

\lablevel{n3}
\labnumber{Laboratorio 1}
\labtitle{Introducción a las listas}

\makeheader

\begin{document}

\maketitle

\begin{sectionbox}{Preparación del ambiente de trabajo}
Cree un nuevo módulo de funciones llamado \texttt{listas.py} en el que escribirá las funciones correspondientes a las actividades 1 a 3 de este laboratorio.
\end{sectionbox}


\begin{sectionbox}{Métodos de listas}

A continuación se presentan algunos métodos de listas útiles para este laboratorio (y para el curso):

\begin{itemize}
    \item \texttt{append(elemento)}: agrega \texttt{elemento} al final de la lista. E.g., \texttt{listica.append(\textquotedbl cosa\textquotedbl)} cuando \texttt{listica} es \texttt{[\textquotedbl elemento\textquotedbl, \textquotedbl cosita\textquotedbl]} convierte a \texttt{listica} en \texttt{[\textquotedbl elemento\textquotedbl, \textquotedbl cosita\textquotedbl, \textquotedbl cosa\textquotedbl]}.
    \item \texttt{insert(índice, elemento)}: inserta \texttt{elemento} en la posición indicada por el índice, desplazando los demás elementos a la derecha.
    \item \texttt{remove(elemento)}: elimina la primera aparición de \texttt{elemento} en la lista.
    \item \texttt{pop(índice)}: elimina el elemento en la posición indicada por el índice y lo retorna.
    \item \texttt{sort()}: ordena la lista en el lugar (modifica la lista original; no retorna nada útil).
    \item \texttt{sorted(lista)}: retorna una nueva lista ordenada sin modificar la original.
    \item \texttt{index(elemento)}: retorna la posición de la primera aparición de \texttt{elemento}.
    \item \texttt{cadena.split(separador)}: método de cadenas que divide \texttt{cadena} en una lista de partes usando \texttt{separador} como delimitador. Si no se especifica separador, divide por espacios.
    \item \texttt{separador.join(lista)}: método de cadenas que une los elementos de \texttt{lista} en una sola cadena, insertando \texttt{separador} entre cada par de elementos.
\end{itemize}

Pregúntese: cuáles de estos métodos modifica una lista sin crear una nueva (lo que se denomina \textit{in-place}) y por ende no deben asignarse a una variable cuando se usan, y cuáles de estos métodos retornan una nueva lista y ergo deben asignarse a una variable para tener sentido. 

\end{sectionbox}

\newpage

\begin{sectionbox}{Actividad 1: Corrector de caracteres}
Escriba una función que reciba como parámetros una oración, la posición de una palabra dentro de la oración, la posición de un carácter dentro de esa palabra y un nuevo carácter. La función debe reemplazar el carácter indicado y retornar la oración corregida.

\textbf{Ejemplos:}
\begin{itemize}
    \item \texttt{\textquotedbl Hola Mundx\textquotedbl}, pos. palabra \texttt{1}, pos. carácter \texttt{4}, nuevo \texttt{\textquotesingle o\textquotesingle} → \texttt{\textquotedbl Hola Mundo\textquotedbl}
    \item \texttt{\textquotedbl me llamo Juan\textquotedbl}, pos. palabra \texttt{0}, pos. carácter \texttt{0}, nuevo \texttt{\textquotesingle M\textquotesingle} → \texttt{\textquotedbl Me llamo Juan\textquotedbl}
\end{itemize}
\end{sectionbox}


\begin{sectionbox}{Actividad 2: Rotación de texto delimitado}
Escriba una función que reciba como parámetros una cadena de caracteres, un separador y un entero \texttt{k}. La función debe dividir la cadena usando el separador, rotar los elementos resultantes \texttt{k} posiciones a la derecha y retornar la cadena reconstruida con el mismo separador.

Una rotación a la derecha de \texttt{k} posiciones desplaza los últimos \texttt{k} elementos al inicio de la lista. Presuma que $0 \leq k < $ número de elementos.

\textbf{Ejemplos:}
\begin{itemize}
    \item \texttt{\textquotedbl lunes martes miércoles jueves\textquotedbl}, sep \texttt{\textquotedbl\ \textquotedbl}, \texttt{k=1} → \texttt{\textquotedbl jueves lunes martes miércoles\textquotedbl}
    \item \texttt{\textquotedbl a-b-c-d-e\textquotedbl}, sep \texttt{\textquotedbl-\textquotedbl}, \texttt{k=0} → \texttt{\textquotedbl a-b-c-d-e\textquotedbl}
\end{itemize}
\end{sectionbox}


\begin{sectionbox}{Actividad 3: Actualizar una playlist}
Escriba una función que reciba como parámetros una lista de canciones, el nombre de una canción a eliminar y el nombre de una canción a agregar. La función debe:

\begin{itemize}
    \item Eliminar la primera aparición de la canción indicada si esta se encuentra en la lista. Si no está, no hacer ninguna modificación en ese paso.
    \item Agregar la nueva canción al final de la lista.
    \item Retornar la lista ordenada alfabéticamente.
\end{itemize}

\textbf{Ejemplos:}
\begin{itemize}
    \item Playlist: \texttt{[\textquotedbl Bohemian Rhapsody\textquotedbl, \textquotedbl Hotel California\textquotedbl, \textquotedbl Stairway to Heaven\textquotedbl]}, eliminar \texttt{\textquotedbl Hotel California\textquotedbl}, agregar \texttt{\textquotedbl Imagine\textquotedbl} → \texttt{[\textquotedbl Bohemian Rhapsody\textquotedbl, \textquotedbl Imagine\textquotedbl, \textquotedbl Stairway to Heaven\textquotedbl]}
    \item \texttt{[\textquotedbl Shape of You\textquotedbl, \textquotedbl Blinding Lights\textquotedbl]}, eliminar \texttt{\textquotedbl Watermelon Sugar\textquotedbl} (no está en la lista), agregar \texttt{\textquotedbl Levitating\textquotedbl} → \texttt{[\textquotedbl Blinding Lights\textquotedbl, \textquotedbl Levitating\textquotedbl, \textquotedbl Shape of You\textquotedbl]}
\end{itemize}
\end{sectionbox}


\begin{sectionbox}{Entrega}
Entregue el archivo \texttt{listas.py} a través de Bloque Neón en el laboratorio del Nivel 3 designado como ``N3-L1: Introducción a las listas''.
\end{sectionbox}

\end{document}
